\makeatletter
\def\input@path{{./}{../}{../../}{../../../}{../../../../}{preamble/}{../preamble/}{../../preamble/}{../../../preamble/}{../../../../preamble/}}
\makeatother
\input{preamble_exams.tex}
\input{graphs/G11_Histograms.tex}

\title{Grade 11 -- Term 1 \\ C5, C6 and C7 Exam Solutions}
\author{}
\date{}

\begin{document}
\maketitle
\vspace{-2cm}
\setcounter{tocdepth}{1}
\setcounter{secnumdepth}{2}
\phantomsection\label{toc}
\tableofcontents

\section*{Learning objective}\vspace{-5mm}
Use continuous probability models and sample data to describe typical values,
variability, and interval probabilities, and interpret differences between
sample estimates and theoretical values.

\section*{Criteria assessment}\vspace{-5mm}
This exam assesses C5, C6 and C7 separately, with 6 points per criterion
(18 points in total). Each task shows its point allocation. Partial credit
may be awarded when the correct procedure is clearly shown but arithmetic
errors are made. Answers without the required procedure will not receive credit.
Failure to write your name on the task sheet or submitting it after the deadline
will result in a one-point deduction from each criterion.
For this exam, a minimum of 4 out of 6 points is required to pass each criterion.
Use unrounded intermediate values; round final statistics to three decimal
places and percentage errors to two decimal places where necessary.

\section*{Evaluated criteria}
\begin{enumerate}[label=\textbf{C\arabic*},start=5]
\item Calculates and interprets the mean, variance and standard deviation of continuous distributions using standard formulas and functions.
\item Estimates interval probabilities using relative frequencies from generated samples.
 and compares them with theoretical probabilities.
\item Calculates sample means, variances, and standard deviations for different sample sizes and data representations.
\end{enumerate}


\sectionScored{Two packing stations}{C5-\sml{2}}{sec:c5-1}
{
A gift company uses two packing stations. At station A, an automated machine
finishes a gift box in anywhere from 4 to 16 minutes. Within these limits,
any two time windows of the same length are equally likely to contain the
finishing time. Let $A$ be its packing time in minutes.

At station B, staff pack custom gift boxes, also taking between 4 and 16
minutes. To describe many orders, the manager draws a smooth outline of how
often different packing times occur. The outline starts at zero at 4 minutes,
rises along a straight edge to its single highest point at 7 minutes, then
falls along a straight edge to zero at 16 minutes. No packing times occur
outside these limits. Let $B$ be this packing time in minutes.
\QuestionsTasks[sec:c5-1]{
{5}/{sol:c5-1a}/{Identify the distribution for A from the clues and calculate its mean, variance, and standard deviation using the corresponding formulas. \score{1}},
{5}/{sol:c5-1b}/{Identify the distribution for B from the clues and calculate its mean, variance, and standard deviation using the corresponding formulas. \score{1}}
}
\phantomsection\label{sol:c5-1a}
\subsubsection*{C5.1(a) -- Uniform model \score{1}}
Equal-length time windows being equally likely, with no times outside
4 to 16 minutes, identifies a uniform distribution: $A\sim U(4,16)$.
With $a=4$ and $b=16$, use $\mu=(a+b)/2$,
$\sigma^2=(b-a)^2/12$, and $\sigma=\sqrt{\sigma^2}$:
\[
E(A)=\frac{4+16}{2}=10\text{ min},\qquad
Var(A)=\frac{(16-4)^2}{12}=12\text{ min}^2,
\qquad \sigma_A=\sqrt{12}\approx3.464\text{ min}.
\]
\phantomsection\label{sol:c5-1b}
\subsubsection*{C5.1(b) -- Triangular model \score{1}}
The two straight edges meeting at one peak identify a triangular distribution.
The endpoints give the minimum $a=4$ and maximum $b=16$; the peak gives
the mode $c=7$. Thus $B\sim\operatorname{Tri}(4,16,7)$, using the order
(minimum, maximum, mode). Apply $\mu=(a+b+c)/3$ and
$\sigma^2=(a^2+b^2+c^2-ab-ac-bc)/18$, then take the square root:
\[
E(B)=\frac{4+16+7}{3}=9\text{ min},
\]
\[
Var(B)=\frac{4^2+16^2+7^2-4(16)-4(7)-16(7)}{18}
=\frac{321-204}{18}=\frac{13}{2}=6.5\text{ min}^2.
\]
\[
\sigma_B=\sqrt{6.5}\approx2.55\text{ min}.
\]
}

\sectionScored{Two custom-order workshops}{C5-\sml{2}}{sec:c5-2}
{
Two workshops make custom wooden signs. Each order takes between 2 and
11 hours to complete. The owners use smooth outlines to describe how often
different completion times occur over many orders; times outside these limits
never occur.

At workshop I, longer jobs are more common: the outline is a single straight
edge rising from zero at 2 hours to its highest point at 11 hours.
At workshop D, shorter jobs are more common: the outline is a single straight
edge falling from its highest point at 2 hours to zero at 11 hours.
Let $I$ and $D$ be the respective completion times in hours.
\QuestionsTasks[sec:c5-2]{
{5}/{sol:c5-2a}/{Identify each distribution from the clues and calculate its mean, variance, and standard deviation using the corresponding formulas. \score{2}},
}
\phantomsection\label{sol:c5-2a}
\subsubsection*{C5.2(a) -- Linear models \score{2}}
For I, the straight rise from zero at the lower endpoint to a peak at the
upper endpoint identifies a linearly increasing distribution. For D, the
straight fall from a peak at the lower endpoint to zero at the upper endpoint
identifies a linearly decreasing distribution. Both have $a=2$ and $b=11$:
$I\sim\operatorname{Lin}_{\mathrm{inc}}(2,11)$ and
$D\sim\operatorname{Lin}_{\mathrm{dec}}(2,11)$. Apply the corresponding mean formulas:
\[
E(I)=\frac{a+2b}{3}=\frac{2+2(11)}{3}=8\text{ h},\qquad
E(D)=\frac{2a+b}{3}=\frac{2(2)+11}{3}=5\text{ h}.
\]
For either linear model,
\[
Var(I)=Var(D)=\frac{(b-a)^2}{18}=\frac{9^2}{18}=4.5\text{ h}^2,
\qquad \sigma_I=\sigma_D=\sqrt{4.5}\approx2.121\text{ h}.
\]

}
\sectionScored{Arrivals and daily revenue}{C5-\sml{2}}{sec:c5-3}
{
A bicycle repair shop receives online service requests throughout its opening
hours, averaging 8 requests per hour. Requests arrive independently, one at
a time, and the arrival pace stays the same throughout the day. In any very
short time window, the chance of a request is proportional to the window's
length, while two or more requests in that window are negligible. How long
the shop has already waited does not affect the chance of a request during
the next minute. Let $T$ be the time between consecutive requests in minutes.

On comparable working days, the shop's total revenue varies because of many
small, independent influences, with none dominating. A smooth outline of
daily revenues has a symmetric bell shape, highest at \$2400 and tapering
on both sides. The two places where its sides change from becoming steeper
to becoming flatter are \$2250 and \$2550. Treat this bell-shaped pattern
as exact for the calculations. Let $R$ be daily revenue measured in hundreds
of dollars.
\QuestionsTasks[sec:c5-3]{
{5}/{sol:c5-3a}/{Identify the distribution for T from the clues and use its formulas to find the mean, variance, and standard deviation in minutes. \score{1}},
{5}/{sol:c5-3b}/{Identify the revenue distribution from the clues. Use its formulas to find the mean, variance, and standard deviation in dollars. Interpret the mean. \score{1}}
}
\phantomsection\label{sol:c5-3a}
\subsubsection*{C5.3(a) -- Exponential waiting time \score{1}}
Independent, single arrivals at a steady pace, together with the fact that
past waiting does not affect future waiting, identify an exponential
distribution for the time between requests. Convert the hourly arrival rate
to minutes: $\lambda=8/60=2/15$ requests per minute.
Thus $T\sim\operatorname{Exp}(2/15)$. Apply the waiting-time formulas:
\[
E(T)=\frac{1}{\lambda}=\frac{15}{2}=7.5\text{ min},\qquad
Var(T)=\frac{1}{\lambda^2}=\frac{225}{4}=56.25\text{ min}^2,
\qquad \sigma_T=\frac{1}{\lambda}=7.5\text{ min}.
\]
\phantomsection\label{sol:c5-3b}
\subsubsection*{C5.3(b) -- Normal revenue and units \score{1}}
The symmetric bell shape and many small independent influences identify a
normal distribution. Its center is its mean, \$2400. For a normal curve,
the changes in curvature occur one standard deviation on either side of
the mean, so $\sigma_Y=2550-2400=2400-2250=150$ dollars.
In hundreds of dollars, $\mu_R=24$ and $\sigma_R=1.5$, giving
$\sigma_R^2=(1.5)^2=2.25$ and $R\sim N(24,2.25)$ (mean, variance).
Let $Y=100R$ be revenue in dollars. Apply
$E(Y)=100\mu_R$, $Var(Y)=100^2\sigma_R^2$, and
$\sigma_Y=\sqrt{Var(Y)}$:
\[
E(Y)=100(24)=2400\text{ dollars},\qquad
Var(Y)=100^2(2.25)=22500\text{ dollars}^2,
\]
\[
\sigma_Y=\sqrt{22500}=150\text{ dollars}.
\]
The model predicts a long-run average daily revenue of \$2400; individual
days need not have that revenue.
\secInicio[sec:c5-3]
}

\sectionScored{Water consumption, raw data}{C6-\sml{1}}{sec:c6-1}
{
A simulated sample of daily household water consumption (hundreds of liters) is
\[
\begin{array}{cccccccccc}
2.1&2.5&2.9&3.0&3.2&3.6&3.9&4.0&4.3&4.7\\
2.4&2.8&3.1&3.4&3.7&4.0&4.1&4.5&4.8&5.2
\end{array}
\]
The theoretical probability is $P(3\le X\le4)=0.5$.
\QuestionsTasks[sec:c6-1]{
{6}/{sol:c6-1a}/{Find the empirical probability of the stated interval from the given sample. \score{0.5}},
{6}/{sol:c6-1b}/{Calculate the absolute and percentage errors relative to the given theoretical probability. \score{0.5}}
}
\phantomsection\label{sol:c6-1a}
\subsubsection*{C6.1(a) -- Counting observations \score{0.5}}
There are $n=20$ observations. Those satisfying both inclusive endpoints are
\[
3.0,\ 3.2,\ 3.6,\ 3.9,\ 4.0,\ 3.1,\ 3.4,\ 3.7,\ 4.0.
\]
Thus $k=9$ and $\widehat P(3\le X\le4)=9/20=0.45$.
\phantomsection\label{sol:c6-1b}
\subsubsection*{C6.1(b) -- Errors \score{0.5}}
\[
E=|0.45-0.5|=0.05,\qquad
\text{Percentage error}=\frac{0.05}{0.5}\times100\%=10\%.
\]
The empirical result is 0.05, or 5 percentage points, below the theoretical
probability. The relative percentage error is 10\%.
\secInicio[sec:c6-1]
}

\sectionScored{Service times, frequency data}{C6-\sml{1}}{sec:c6-2}
{
Recorded service times are summarized below.
\[
\begin{array}{c|rrrrrr}
 x_i\text{ (min)}&3&6&9&12&15&18\\\hline
 f_i&6&18&30&24&15&7
\end{array}
\]
The theoretical probability is $P(6\le X<15)=0.75$.
\QuestionsTasks[sec:c6-2]{
{6}/{sol:c6-2a}/{Find the empirical probability of the stated interval from the given sample. \score{0.5}},
{6}/{sol:c6-2b}/{Calculate the absolute and percentage errors relative to the given theoretical probability. \score{0.5}}
}
\phantomsection\label{sol:c6-2a}
\subsubsection*{C6.2(a) -- Frequency totals \score{0.5}}
\[
n=6+18+30+24+15+7=100,\qquad k=18+30+24=72.
\]
The value 15 is excluded, so $\widehat P(6\le X<15)=72/100=0.72$.
\phantomsection\label{sol:c6-2b}
\subsubsection*{C6.2(b) -- Errors \score{0.5}}
\[
E=|0.72-0.75|=0.03,\qquad
\text{Percentage error}=\frac{0.03}{0.75}\times100\%=4\%.
\]
\secInicio[sec:c6-2]
}

\sectionScored{Download times, grouped data}{C6-\sml{2}}{sec:c6-3}
{
A sample of download times is summarized by the following frequencies.
All times are in seconds; classes are left-inclusive and right-exclusive.
\[
\begin{array}{c|rrrr}
\text{Class}&[0,5)&[5,10)&[10,15)&[15,20)\\\hline
f_i&4&8&6&2
\end{array}
\]
The model gives $P(5\le X<15)=0.8$.
\QuestionsTasks[sec:c6-3]{
{6}/{sol:c6-3a}/{Find the empirical probability of the stated interval from the given sample. \score{1}},
{6}/{sol:c6-3b}/{Calculate the absolute and percentage errors relative to the given theoretical probability. \score{1}}
}
\phantomsection\label{sol:c6-3a}
\subsubsection*{C6.3(a) -- Grouped frequencies \score{1}}
The event contains the two middle classes, so
\[
n=4+8+6+2=20,\qquad k=8+6=14,\qquad
\widehat P(5\le X<15)=\frac{14}{20}=0.7.
\]
The event is a union of complete classes, so its relative frequency can
be calculated without estimating individual download times.
\phantomsection\label{sol:c6-3b}
\subsubsection*{C6.3(b) -- Errors \score{1}}
\[
E=|0.7-0.8|=0.1,\qquad
\text{Percentage error}=\frac{0.1}{0.8}\times100\%=12.5\%.
\]
\secInicio[sec:c6-3]
}

\sectionScored{Ticket queues, histogram data}{C6-\sml{2}}{sec:c6-4}
{
The histogram is the only record of a sample of queue times. Bar heights
are frequencies; times are in minutes, with left-inclusive, right-exclusive classes.
\begin{center}
\HistogramGraph{(0,3),(3,6),(6,9),(9,12)}{6,12,9,3}
{0.60\textwidth}{4.5cm}{blue!25}{3}{3}
\end{center}
The theoretical probability is $P(3\le X<9)=0.75$.
\QuestionsTasks[sec:c6-4]{
{6}/{sol:c6-4a}/{Find the empirical probability of the stated interval from the given sample. \score{1}},
{6}/{sol:c6-4b}/{Calculate the absolute and percentage errors relative to the given theoretical probability. \score{1}}
}
\phantomsection\label{sol:c6-4a}
\subsubsection*{C6.4(a) -- Reading the histogram \score{1}}
The classes $[0,3),[3,6),[6,9),[9,12)$ have respective frequencies
$6,12,9,3$. Hence $n=6+12+9+3=30$ and $k=12+9=21$.
\[
\widehat P(3\le X<9)=\frac{21}{30}=0.7.
\]
\phantomsection\label{sol:c6-4b}
\subsubsection*{C6.4(b) -- Errors \score{1}}
\[
E=|0.7-0.75|=0.05,\qquad
\text{Percentage error}=\frac{0.05}{0.75}\times100\%
=\frac{20}{3}\%\approx6.67\%.
\]
\secInicio[sec:c6-4]
}

Use the sample-variance denominator $n-1$ throughout these four problems.
For each statistic, calculate absolute error $E=|\text{sample value}-\text{theoretical value}|$
and percentage error $100E/|\text{theoretical value}|$. Use unrounded values in all error calculations.
For grouped data and histograms, use class midpoints to estimate the statistics.

\sectionScored{Sensor readings, raw data}{C7-\sml{1}}{sec:c7-1}
{
Five sensor readings in degrees Celsius are $9,11,12,13,15$.
The theoretical model has $\mu=12$ degrees Celsius and
$\sigma^2=4$ square degrees Celsius.
\QuestionsTasks[sec:c7-1]{
{7}/{sol:c7-1a}/{Calculate the sample mean, sample variance, and sample standard deviation from the given data. \score{0.5}},
{7}/{sol:c7-1b}/{Calculate the absolute and percentage errors of the mean, variance, and standard deviation relative to the theoretical values. \score{0.5}}
}
\phantomsection\label{sol:c7-1a}
\subsubsection*{C7.1(a) -- Raw-data statistics \score{0.5}}
\[
 n=5,\qquad\bar{x}=\frac{9+11+12+13+15}{5}=\frac{60}{5}=12.
\]
\[
\begin{array}{c|rrrrr}
x_i&9&11&12&13&15\\\hline
x_i-\bar{x}&-3&-1&0&1&3\\
(x_i-\bar{x})^2&9&1&0&1&9
\end{array}
\]
The squared deviations total $20$. Thus
\[
 s^2=\frac{20}{5-1}=5,\qquad s=\sqrt{5}\approx2.236.
\]
The mean and standard deviation are in degrees Celsius; the variance is
in square degrees Celsius.
\phantomsection\label{sol:c7-1b}
\subsubsection*{C7.1(b) -- Absolute and percentage errors \score{0.5}}
The theoretical standard deviation is $\sigma=\sqrt{\sigma^2}$.
Apply the error formulas to each statistic using the unrounded results above:
\[
\begin{array}{l|rr}
\text{Statistic}&\text{Absolute error}&\text{Percentage error}\\\hline
Mean&0\phantom{.000}&0\phantom{.00}\%\\
Variance&1\phantom{.000}&25\phantom{.00}\%\\
Standard deviation&0.236&11.8\phantom{0}\%\\
\end{array}
\]
Absolute errors have the units of the corresponding statistic; percentage errors are unitless.
\secInicio[sec:c7-1]

}
\sectionScored{Items sold, frequency data}{C7-\sml{1}}{sec:c7-2}
{
The numbers of items sold in six short sales periods are summarized below.
The theoretical model has $\mu=5$ items and $\sigma^2=5$ items$^2$.
\[
\begin{array}{c|rrr}x_i&2&5&8\\\hline f_i&1&3&2\end{array}
\]
\QuestionsTasks[sec:c7-2]{
{7}/{sol:c7-2a}/{Calculate the sample mean, sample variance, and sample standard deviation from the given data. \score{0.5}},
{7}/{sol:c7-2b}/{Calculate the absolute and percentage errors of the mean, variance, and standard deviation relative to the theoretical values. \score{0.5}}
}
\phantomsection\label{sol:c7-2a}
\subsubsection*{C7.2(a) -- Weighted statistics \score{0.5}}
\[
 n=1+3+2=6,\qquad\bar{x}=\frac{1(2)+3(5)+2(8)}{6}
 =\frac{33}{6}=5.5\text{ items}.
\]
The weighted squared deviations are
\[
1(2-5.5)^2=12.25,\quad 3(5-5.5)^2=0.75,\quad
2(8-5.5)^2=12.5.
\]
\[
s^2=\frac{12.25+0.75+12.5}{6-1}=\frac{25.5}{5}=5.1\text{ items}^2,
\qquad s=\sqrt{5.1}\approx2.258\text{ items}.
\]
\phantomsection\label{sol:c7-2b}
\subsubsection*{C7.2(b) -- Absolute and percentage errors \score{0.5}}
The theoretical standard deviation is $\sigma=\sqrt{\sigma^2}$.
Apply the error formulas to each statistic using the unrounded results above:
\[
\begin{array}{l|rr}
\text{Statistic}&\text{Absolute error}&\text{Percentage error}\\\hline
Mean&0.5\phantom{00}&10\phantom{.00}\%\\
Variance&0.1\phantom{00}&2\phantom{.00}\%\\
Standard deviation&0.022&1\phantom{.00}\%\\
\end{array}
\]
Absolute errors have the units of the corresponding statistic; percentage errors are unitless.
\secInicio[sec:c7-2]
}

\sectionScored{Cycling times, grouped data}{C7-\sml{2}}{sec:c7-3}
{
Cycling times in minutes are grouped as follows, using left-inclusive,
right-exclusive classes. The theoretical mean is 18 minutes and variance
is 20 minutes$^2$.
\[
\begin{array}{c|rrr}
\text{Time (min)}&[10,14)&[14,18)&[18,22)\\\hline
f_i&1&2&3
\end{array}
\]
\QuestionsTasks[sec:c7-3]{
{7}/{sol:c7-3a}/{Calculate the sample mean, sample variance, and sample standard deviation from the given data. \score{1}},
{7}/{sol:c7-3b}/{Calculate the absolute and percentage errors of the mean, variance, and standard deviation relative to the theoretical values. \score{1}}
}
\phantomsection\label{sol:c7-3a}
\subsubsection*{C7.3(a) -- Midpoint estimates \score{1}}
The midpoints are $(10+14)/2=12$, $(14+18)/2=16$, and
$(18+22)/2=20$. The sample size is $n=1+2+3=6$.
\[
\bar{x}\approx\frac{1(12)+2(16)+3(20)}{6}
=\frac{104}{6}=\frac{52}{3}\approx17.333\text{ min}.
\]
Using the unrounded midpoint mean, the squared-deviation sum is
\[
1\left(12-\frac{52}{3}\right)^2+
2\left(16-\frac{52}{3}\right)^2+
3\left(20-\frac{52}{3}\right)^2
=\frac{256+32+192}{9}=\frac{160}{3}.
\]
\[
s^2\approx\frac{160/3}{6-1}=\frac{32}{3}\approx10.667\text{ min}^2,
\qquad s\approx\sqrt{\frac{32}{3}}\approx3.266\text{ min}.
\]
\phantomsection\label{sol:c7-3b}
\subsubsection*{C7.3(b) -- Absolute and percentage errors \score{1}}
The theoretical standard deviation is $\sigma=\sqrt{\sigma^2}$.
Apply the error formulas to each statistic using the unrounded results above:
\[
\begin{array}{l|rr}
\text{Statistic}&\text{Absolute error}&\text{Percentage error}\\\hline
Mean&0.667&3.7\phantom{0}\%\\
Variance&9.333&46.67\%\\
Standard deviation&1.206&26.97\%\\
\end{array}
\]
Absolute errors have the units of the corresponding statistic; percentage errors are unitless.
These errors use midpoint estimates because individual observations are unknown.
\secInicio[sec:c7-3]
}

\sectionScored{Crate masses, histogram data}{C7-\sml{2}}{sec:c7-4}
{
A histogram gives a sample of crate masses in kilograms. Bar heights are
frequencies; classes are left-inclusive and right-exclusive. The model
has $\mu=5$ kg and $\sigma^2=2$ kg$^2$.
\begin{center}
\HistogramGraph{(2,4),(4,6),(6,8)}{2,3,1}
{0.55\textwidth}{4.5cm}{orange!30}{2}{1}
\end{center}
\QuestionsTasks[sec:c7-4]{
{7}/{sol:c7-4a}/{Calculate the sample mean, sample variance, and sample standard deviation from the given data. \score{1}},
{7}/{sol:c7-4b}/{Calculate the absolute and percentage errors of the mean, variance, and standard deviation relative to the theoretical values. \score{1}}
}
\phantomsection\label{sol:c7-4a}
\subsubsection*{C7.4(a) -- Histogram statistics \score{1}}
The classes $[2,4),[4,6),[6,8)$ have frequencies $2,3,1$, respectively.
Thus $n=6$; the midpoints are $3,5,7$ kg.
\[
\bar{x}\approx\frac{2(3)+3(5)+1(7)}{6}=\frac{28}{6}=\frac{14}{3}
\approx4.667\text{ kg}.
\]
The weighted squared-deviation sum is
\[
2\left(3-\frac{14}{3}\right)^2+
3\left(5-\frac{14}{3}\right)^2+
\left(7-\frac{14}{3}\right)^2
=\frac{50+3+49}{9}=\frac{34}{3}.
\]
\[
s^2\approx\frac{34/3}{6-1}=\frac{34}{15}\approx2.267\text{ kg}^2,
\qquad s\approx\sqrt{\frac{34}{15}}\approx1.506\text{ kg}.
\]
\phantomsection\label{sol:c7-4b}
\subsubsection*{C7.4(b) -- Absolute and percentage errors \score{1}}
The theoretical standard deviation is $\sigma=\sqrt{\sigma^2}$.
Apply the error formulas to each statistic using the unrounded results above:
\[
\begin{array}{l|rr}
\text{Statistic}&\text{Absolute error}&\text{Percentage error}\\\hline
Mean&0.333&6.67\%\\
Variance&0.267&13.33\%\\
Standard deviation&0.091&6.46\%\\
\end{array}
\]
Absolute errors have the units of the corresponding statistic; percentage errors are unitless.
These errors use midpoint estimates because individual observations are unknown.
\secInicio[sec:c7-4]

\newpage
\section{Summarized Solutions}
\subsection*{C5}
\begin{enumerate}[label=\textbf{C5.\arabic*.}]
\item (a) $A\sim U(4,16)$; $E(A)=10$ min, $Var(A)=12$ min$^2$, $\sigma_A\approx3.464$ min.
(b) $B\sim\operatorname{Tri}(4,16,7)$ (minimum, maximum, mode);
$E(B)=9$ min, $Var(B)=6.5$ min$^2$, $\sigma_B\approx2.55$ min.
\item (a) I is linearly increasing and D is linearly decreasing on $[2,11]$.
$E(I)=8$ h, $E(D)=5$ h; both variances $4.5$ h$^2$ and
standard deviations $\approx2.121$ h.
\item (a) $T\sim\operatorname{Exp}(2/15)$ in minutes;
$E(T)=7.5$ min, $Var(T)=56.25$ min$^2$, $\sigma_T=7.5$ min.
(b) $R\sim N(24,2.25)$ (mean, variance) in hundreds of dollars.
Dollar revenue: mean \$2400, variance $22500$ dollars$^2$,
standard deviation \$150; long-run daily average \$2400.
\end{enumerate}
\subsection*{C6}
\begin{enumerate}[label=\textbf{C6.\arabic*.}]
\item (a) $n=20$, $k=9$, $\widehat P=0.45$.
(b) Absolute error $0.05$; percentage error $10\%$.
\item (a) $n=100$, $k=72$, $\widehat P=0.72$.
(b) Absolute error $0.03$; percentage error $4\%$.
\item (a) $n=20$, $k=14$, $\widehat P=0.7$.
(b) Absolute error $0.1$; percentage error $12.5\%$.
\item (a) Frequencies $(6,12,9,3)$, $n=30$, $k=21$, $\widehat P=0.7$.
(b) Absolute error $0.05$; percentage error $\approx6.67\%$.
\end{enumerate}
\subsection*{C7}
\begin{enumerate}[label=\textbf{C7.\arabic*.}]
\item (a) $n=5$, $\bar{x}=12$, $s^2=5$, $s\approx2.236$
(mean and standard deviation in degrees Celsius, variance in square degrees).
(b) Absolute errors (mean, variance, SD): $(0,1,0.236)$; percentage errors: $(0\%,25\%,11.8\%)$.
\item (a) $n=6$, $\bar{x}=5.5$ items, $s^2=5.1$ items$^2$,
$s\approx2.258$ items.
(b) Absolute errors (mean, variance, SD): $(0.5,0.1,0.022)$; percentage errors: $(10\%,2\%,1\%)$.
\item (a) Midpoints $(12,16,20)$ min; $n=6$; midpoint estimates
$\bar{x}\approx17.333$ min, $s^2\approx10.667$ min$^2$,
$s\approx3.266$ min.
(b) Absolute errors (mean, variance, SD): $(0.667,9.333,1.206)$; percentage errors: $(3.7\%,46.67\%,26.97\%)$.
\item (a) Frequencies $(2,3,1)$, $n=6$, midpoints $(3,5,7)$ kg;
midpoint estimates $\bar{x}\approx4.667$ kg, $s^2\approx2.267$ kg$^2$,
$s\approx1.506$ kg.
(b) Absolute errors (mean, variance, SD): $(0.333,0.267,0.091)$; percentage errors: $(6.67\%,13.33\%,6.46\%)$.
\end{enumerate}
}

\end{document}
